\answer{ex:15:1}{(a)~$2\cdot10^5$ per neuron, $6\cdot10^{12}$ for the circuit. (b)~$\approx8$~bit/s per wire, at least $2.5\cdot10^4$~s $\approx7$~h.}
\answer{ex:15:2}{Factor $2$: neighbor correlation $0.943$, $\lambda$ from $4.18$ to $7.3\cdot10^{-4}$, ratio $\approx5.7\cdot10^3$; factor $4$: correlation $0.8$, ratio $\approx94$.}
\answer{ex:15:3}{(a)~$\lambda=0.988$ and $0.808$: $82$ and $4.7$ movements. (b)~$x_s^*=0.690$, $x_f^*=0.172$, sum $0.862$.}
\answer{ex:15:4}{(a)~$116$ movements. (b)~$19$. A single-process model with a zero sum gives no savings.}
\answer{ex:15:5}{(a)~$G=50$~s$^{-1}$ versus the limit of $10.5$~s$^{-1}$ without a model. (b)~No: at $G=50$~s$^{-1}$ the critical delay is $d_c\approx103\ms$; at $d=150\ms$ one needs $\hat k>0.445k$ (for any $G$ and $d$, $\hat k>k/2$).}
\answer{ex:15:6}{The error falls below $0.5\%$ in $6$--$30$~s, with a floor of $\approx(6$--$9)\cdot10^{-4}$ against $7.5\cdot10^{-4}$ for the best weights; at $\eta=50$ the weights still differ from the best by up to $0.2$ along the ill-conditioned directions of $C$ (Exercise~\ref{ex:15:2}).}
\answer{ex:15:7}{$\mathbf B=A\mathbf K$; $q_f/R\approx0.035$, $q_s/R\approx8.6\cdot10^{-4}$; with $4q_s$, $B_f\approx0.088$, $B_s\approx0.047$: the slow process learns more than twice as fast, and the fast one more slowly.}
